\section{Conclusion}
A decision that this literature does not report accounts for more of the
variation in its results than the architectural differences that distinguish
the systems being compared, and it changes what every other measurement means.

Holding corpus, architecture, features and seed fixed and varying only what the
test partition may share with training lowers ROC-AUC from $0.796$ to $0.585$,
which is $71.5\%$ of the margin above chance. Holding the protocol fixed instead
and changing only the training recipe, fixed by grid search on the validation
partition with test withheld, raises the repository-disjoint score to $0.650$
and moves that leakage share to $45.5\%$. Both recipes are defensible and
neither is stated in the papers we surveyed, so a leakage figure is a joint
property of a corpus and a training configuration rather than of a benchmark.

The same contingency governs the representation. The reverse-edge encoding gains $+0.033$ under the default recipe and
loses $0.006$ under the selected one, on identical data with identical features;
five-fold averaging behaves the same way, worth $+0.033$ and then $-0.002$. The
relation labels are null under both. A further consequence is that honest
tuning cannot recover the best available number: across the twelve
configurations searched, validation ranks test only weakly ($\rho=0.32$), and
selecting on validation leaves about a point of ROC-AUC on the table. We report
these as conditional measurements rather than as recommendations, because that
is what they are.

Two findings do not depend on the choice. Given the same graphs with node
content replaced by noise, the models recover labels defined by the edge set
alone, gaining $+0.104$ and $+0.303$ ROC-AUC from the edges, so the pipeline
demonstrably reads graph structure. And on the $2{,}067$ vulnerable/patched
pairs inside the test partition, where project identity and coding style cancel
by construction, no configuration separates a function from its own fix at
better than chance: $0.502$ where only structure differs. Aggregate ROC-AUC is
dominated by unrelated safe functions and does not measure that capability.

We present these as measurements rather than as criticism. Published results in
this area appear to be accurately computed under the procedures employed. What
has been missing is a statement of what those procedures are. On the evidence
here, an ablation reported without its training procedure is not interpretable,
and a leakage figure reported without one is not comparable.
